% Fluid Flow vs Euler's Equations vs OpenAI
% Johan Jansson and Claes Johnson. DRAFT -- summing up.
\documentclass[11pt]{article}
\usepackage[margin=1.15in]{geometry}
\usepackage{amsmath,amssymb}
\usepackage{url}
\usepackage[hidelinks]{hyperref}
\setlength{\parskip}{3pt}

\newcommand{\grad}{\nabla}
\newcommand{\eps}{\varepsilon}

\title{\vspace{-2em}Fluid Flow versus Euler's Equations versus OpenAI}
\author{Johan Jansson \qquad Claes Johnson\\[0.4em]
\small KTH Royal Institute of Technology, Stockholm}
\date{\today}

\begin{document}
\maketitle

\begin{abstract}
\noindent
Three things are to be kept apart: a fluid flow, which is a phenomenon; Euler's equations, which are
a model of it; and a proof about the equations, which is a statement about the model. We state the
equations, note that they are parameter free, and observe that the mathematical question asked of
them --- existence, regularity, uniqueness --- has not been answered by analysis in symbolic form,
because the solutions turn out to have a complicated nature which symbolic analysis has no grip on.
Solutions in digital form over a mesh can be computed by a model we name \emph{Adaptive Euler} ---
a stabilised finite element discretisation together with duality-based output error control and
adaptive refinement --- and they are turbulent, fluctuating over a range of scales, with well
determined mean values such as drag and lift which agree with measurement and arrive with a computed
bound on their error. We
then ask what a recent symbolic proof of finite-time blowup establishes, and conclude that it is a
statement about the model which does not reach either of the other two levels: it exhibits no
individual case, so nothing can be inserted into the equations, inspected, or measured.
\end{abstract}

\section{The equations}\label{sec:eq}

Let $\Omega\subset\mathbb{R}^3$ be a domain with boundary $\Gamma$, let $x=(x_1,x_2,x_3)$ be a
Euclidean coordinate in $\Omega$ and $t$ time. An inviscid incompressible fluid of unit density
filling $\Omega$ is described by a velocity $u(x,t)\in\mathbb{R}^3$ and a pressure $p(x,t)$
satisfying, for $0<t\le T$,
\begin{align}
\dot u + (u\cdot\grad)u + \grad p &= f \qquad\text{in }\Omega, \label{eq:mom}\\[2pt]
\grad\cdot u &= 0 \qquad\;\;\text{in }\Omega, \label{eq:inc}
\end{align}
together with the \emph{slip} boundary condition
\begin{equation}\label{eq:slip}
u\cdot n = 0 \qquad\text{on }\Gamma ,
\end{equation}
$n$ the outward normal, and an initial velocity $u(\cdot,0)=u^0$ in $\Omega$.

Equation~\eqref{eq:mom} is Newton's second law. For unit density the mass of a fluid particle is its
volume, so force equals acceleration directly: the left-hand side $\dot u+(u\cdot\grad)u$ is the
acceleration of a particle following the flow, and the right-hand side is the force acting on it ---
the pressure gradient $-\grad p$, which is the force the surrounding fluid exerts, together with any
body force $f$. Equation~\eqref{eq:inc} states incompressibility: the volume of a fluid particle does
not change, so the velocity field is divergence free, and the pressure is whatever field enforces
this. Condition~\eqref{eq:slip} states that the fluid does not pass through the wall, and nothing
more. It is the only boundary condition the equations admit: the normal component is fixed by
impermeability, nothing in~\eqref{eq:mom}--\eqref{eq:inc} determines the tangential velocity at the
wall, and prescribing it as well would over-determine the problem. The no-slip condition belongs to
Navier--Stokes, where a second-order viscous term carries the extra condition; here there is no
viscosity, and therefore no boundary layer either.

\paragraph{What `inviscid' means.} Not that the viscosity is zero. Air and water have viscosity,
and no flow of either is literally inviscid. What is meant is that the Reynolds number is large ---
$10^6$ to $10^8$ in the flows of aerodynamics and ship hydrodynamics --- so that viscous forces are
negligible beside inertial ones everywhere except in layers which are thin and, as it turns out, do
not determine the forces on the body. Euler's equations are the model of that limit, and the limit is
where the flows of practical interest live. The question is never whether a real fluid has viscosity.
It is whether the quantities one wants depend on how much.

\paragraph{The model is parameter free.} There is no constant in
\eqref{eq:mom}--\eqref{eq:slip} to be measured, fitted or chosen. What enters as data is the
\emph{geometry} --- the domain $\Omega$, which is to say the shape of the body the fluid flows past
--- the \emph{forcing} $f$, and the initial state $u^0$. Nothing else. Whatever the equations predict
is therefore a consequence of the geometry and the forcing alone, and any disagreement with
observation is a fault of the model and cannot be repaired by adjustment.

\section{The mathematical question}\label{sec:question}

The question mathematics asks of a system such as
\eqref{eq:mom}--\eqref{eq:slip} is about the \emph{properties of its solutions}. Does a solution
exist for the given data? Is it regular --- does it stay smooth, or do derivatives become unbounded?
Is it unique, so that the data determine the flow?

These are the right questions to ask and they have not been answered in symbolic form. The reason is
not that the problem has been neglected; it is that the solutions turn out to have a complicated
nature. A flow at a Reynolds number of any interest is turbulent: it fluctuates on a range of scales
from the size of the body down to scales many orders smaller, with no symmetry, no similarity
variable and no closed form. Symbolic analysis works by rewriting, and rewriting needs something to
grip. Against an object defined by carrying structure at every scale with no organising principle to
exploit, there is nothing for it to take hold of, and this is not a temporary state of the technique
but a description of what turbulence is.

What symbolic analysis has produced instead are statements about objects which do possess exploitable
structure --- self-similar profiles, symmetric configurations, families built by iteration --- and
conditional bounds of the form: if the velocity is bounded by $A$, then the derivatives are bounded by
an expression exponential in $A$ iterated several times. As theorems these are real. As statements
about a fluid they are empty, since a bound of size $\exp\exp\exp A$ is not one any flow can be said
to respect.

\section{Adaptive Euler: solutions in digital form}\label{sec:digital}

The same equations treated in digital form behave quite differently. We name the computational model
\textbf{Adaptive Euler}. It has three parts and the third is not an afterthought:

\begin{enumerate}
\item the equations \eqref{eq:mom}--\eqref{eq:slip}, discretised by a stabilised finite element
method over a mesh~\cite{book}, the stabilisation acting as a vanishing artificial viscosity so that
what is computed is a \emph{viscosity solution} of Euler, the limit taken as the mesh is refined;
\item \emph{duality-based output error control}: for each output of interest --- a drag, a lift ---
the error is bounded by the residual of the computed solution weighted by the solution of an adjoint
problem, which measures how much the residual matters for that output;
\item \emph{adaptivity}: the mesh is refined where the weighted residual says the error lives, and
not uniformly, until the bound on the output is below what is asked.
\end{enumerate}

The error control is part of the model rather than an assessment of it. What Adaptive Euler returns is
not a flow but a \emph{pair}: an output and a bound on how wrong that output may be, computed
together. A model which returned the first without the second would be a different model and a much
weaker one, since nothing could then be concluded from it.

\paragraph{The viscosity introduced, and why its size does not matter.} A computation at finite
resolution cannot have zero dissipation, so Adaptive Euler necessarily introduces a viscosity. This is
the obvious objection to it --- that what is being solved is not Euler's equations but some nearby
viscous system --- and the answer is a computational fact rather than an argument.

\emph{The mean-value outputs do not depend on the absolute size of the viscosity introduced.} Refine
the mesh: the artificial viscosity falls, the resolved scales multiply, the flow field changes in
every detail --- and the drag does not move. The dissipation rate does not fall with the viscosity
either; it settles at a value fixed by the geometry and the free stream. That independence is the
dissipation anomaly in computational form, the phenomenon Onsager accounted for~\cite{onsager} and
Taylor measured, and it is the reason the limit is meaningful at all.

It also places the computation and the fluid in the \emph{same} position rather than in different
ones. Air has a small viscosity whose absolute size does not determine the drag of a wing; Adaptive
Euler has a small viscosity whose absolute size does not determine the computed drag. Neither is
inviscid and neither needs to be. What both exhibit is a flow in which the viscosity sets the
smallest scale present and nothing that is asked about.

What Adaptive Euler produces has the following properties, and none of them was put in.

\paragraph{It is turbulent.} The computed flow fluctuates, in time and across a range of scales
widening as the mesh is refined. It is not a steady solution, not a symmetric one, and not one that
could have been guessed.

\paragraph{Its mean values are determined.} The flow field itself is not reproducible in detail ---
and could not be, since for weak solutions of Euler uniqueness is false. But mean-value outputs are:
drag, lift, the position of the separation line. These converge under refinement, and each arrives
with the bound of part~2 above, so the answer carries a statement of how wrong it may be.

\paragraph{They agree with measurement.} Drag within a few per cent, the separation line where
experiment puts it, the lift curve and the stall, the drag crisis~\cite{dragcrisis}. These are
quantities measured by other people, and the comparison is not one we administer: the cases, the
geometries and the measurements belong to the benchmark workshops, and predictions are submitted
before the measurements are released. Adaptive Euler has been put through them --- the high-order CFD
workshop case CS1~\cite{cs1}, and validation against the NASA High Lift Prediction
Workshop~4~\cite{ifasd}, both by Jansson --- which is the form of accountability a computational
model can have and an existence proof cannot.

\paragraph{And it is a solution of Euler.} That the limit satisfies \eqref{eq:mom}--\eqref{eq:slip}
weakly is not an identification offered because the numbers look right. The stabilisation of the
method is a weighted least-squares control of the residual; the discrete energy estimate bounds that
residual in a weighted norm; testing against a smooth function sends it to zero in the sense the weak
formulation requires. The limit is a weak solution, as a theorem.

So the questions of Section~\ref{sec:question} do have answers in digital form, though not the
answers the symbolic form was looking for. Existence: by construction and exhibition. Regularity:
none --- the velocity stays bounded while its gradients do not, the solution being H\"older continuous
with exponent about $1/3$ --- the exponent below which Onsager observed that energy conservation can
fail~\cite{onsager} --- and dissipating energy at a rate independent of the vanishing viscosity.
Uniqueness: false for the flow field, true for the mean values, which is the form in which the
question has an affirmative answer and the form in which anybody wants one.

\paragraph{No AI is needed for Euler.} The equations are parameter free, and the parameters of an
AI have no role to play.

That is the whole of it. The equations are given, in symbolic form, in three lines; there is nothing
to infer from data because Newton's second law and incompressibility are not hypotheses to be fitted.
The discretisation is derived from the equations, the error estimate from the residual and the
adjoint, the adaptivity from the estimate. No quantity enters anywhere that was obtained by fitting
to observations. The model carries no parameters, Adaptive Euler carries none either --- a mesh and a
tolerance are choices stated by the user, not constants calibrated against the answer wanted --- and
a large language model carries of the order of $10^{11}$ to $10^{12}$, obtained by optimisation
against a corpus.

Trained closures and surrogates exist and are pursued, and we make no claim about them in general.
We observe only that any such component reintroduces fitted parameters, and that this is the property
Section~\ref{sec:eq} found the model to be free of. A parameter-free prediction which disagrees with
measurement condemns the model. A prediction carrying a trained component which disagrees invites
retraining. The first is a scientific statement and the second is a procedure.

\section{Inspection of blow-up proofs by OpenAI}\label{sec:openai}

Against this stands a recent proof, produced by an AI system and formalised in Lean, that there
exists smooth, compactly supported, divergence-free initial data on $\mathbb{R}^3$ whose smooth Euler
solution has finite maximal lifespan: the velocity gradient becomes unbounded, and the
Beale--Kato--Majda~\cite{bkm} vorticity integral diverges~\cite{openaieuler}. Comparable results were being
prepared independently and by hand~\cite{buckmaster}. We take the proof to be correct, and nothing
below depends on whether it is.

It is a statement about the \emph{model}, and it is a real one: a question open for decades is
settled. Three observations place it.

\paragraph{The velocity does not become infinite.} What diverges is the gradient. That is the same
signature as the computed turbulent solution --- gradients unbounded, velocity bounded --- so the two
accounts agree that smooth solutions of \eqref{eq:mom}--\eqref{eq:slip} break down in finite time,
which is what we reported computationally in~\cite{blowupHJ}. The disagreement lies elsewhere.

\paragraph{It conserves energy, and therefore cannot be a fluid.} With smooth data the solution is
classical up to the blowup time, and a classical solution of \eqref{eq:mom}--\eqref{eq:inc} with
decaying data satisfies
\[
\frac{d}{dt}\int \tfrac12 |u|^2 = 0 ,
\]
by an integration by parts. The construction therefore dissipates \emph{nothing} over the whole of
its existence. The singularity is not a mechanism by which energy is lost but one by which it is
piled up: at the singular time all of it is still present, gathered into a vanishing region. A fluid
loses energy continuously from the moment it meets a body, and the loss is the drag.

\paragraph{It exhibits no individual case.} The datum is the limit of an iteration over stages and
has no closed form. Nothing can be plotted, nothing measured. Nor can either check on a \emph{result}
be applied: there is no symbolic expression to insert into
\eqref{eq:mom}--\eqref{eq:inc} and reduce, and no numerical field either, since no finite stage of
the iteration blows up --- the property belongs to the limit. What exists is the existence.

\paragraph{What is proved was already claimed.} There is a further observation, and it is not
ungenerous. The theorem asserts that \emph{some} smooth datum has finite lifespan. Our position,
stated in~\cite{blowupHJ} in 2008 and resting on computation, is stronger: \emph{every} laminar Euler
flow at high Reynolds number is unstable and breaks down, with the velocity staying bounded while the
gradients do not. The existence statement follows from ours as a corollary; if all smooth solutions
break down then certainly one does.

What is new is therefore the proof and not the fact. That is a genuine contribution --- a fact
asserted on computational evidence is not the same as a fact established --- and we do not diminish
it. But it should be seen for what it is: a demonstration, in a special constructed case, of a
phenomenon which is generic, which has been exhibited computationally for two decades, and which the
instability of Section~\ref{sec:question} predicts of every flow past every body.

\paragraph{And the viscous case delivers less than it appears to.} The companion proof for
Navier--Stokes constructs a singularity at prescribed forcing, and the Navier--Stokes equations
rescale: if $v$ solves them at unit viscosity then $u(x,t)=a\,v(bx,ct)$ solves them at viscosity
$\nu$ for $a=\nu b$ and $c=ab$, with forcing $a^2b\,g$. The family is therefore generated by a single
unit-viscosity construction, and one must choose how to run it down to the viscosity of a fluid. Hold
the amplitude at $a=1$ and the forcing required grows as $\nu^{-1}$, which leaves the admissible
data. Hold the forcing at $a^2b=1$ and the amplitude, the length and the time all go to zero --- at
$\nu=10^{-6}$ the velocity reaches $10^{-2}$ on a scale of $10^{-4}$ before blowing up. The
singularity shrinks to nothing in the limit in which fluid mechanics lies.

So the construction requires either a force no fluid is subject to or a singularity too small to
constitute anything, and Horace supplies the description: \emph{parturient montes, nascetur
ridiculus mus}. The machinery is formidable and what emerges from it is small.

\section{Three levels, and what reaches which}\label{sec:levels}

The title names three things and the point of the paper is that they are three.

\emph{A fluid flow} is a phenomenon. It has a drag which can be measured.

\emph{Euler's equations} are a model of it, parameter free, with geometry and forcing as data. They
are answerable to the phenomenon: if the drag they give is wrong, nothing can be adjusted to repair
it.

\emph{A proof about the equations} is a statement about the model. It is answerable to the model
alone.

The digital solution reaches all three levels. It is a solution of the equations, by the residual
control; it is an object that can be inspected, and inspecting such objects is where the resolution
of d'Alembert's paradox~\cite{dalembert} and the account of flight~\cite{flight} came from; and its
mean values can be set beside a measurement, where they survive.

The symbolic proof reaches the first of the three and stops. That is not a criticism of it as
mathematics, where it is an advance, and existence without an exhibited instance is ordinary and
accepted practice --- the disproof of the Mertens conjecture~\cite{mertens} established that a
counterexample exists and exhibited none. It is a statement about what the result can be offered as.
A theorem about a model becomes a statement about fluids only by way of an instance, and there is
none here.

\section{The shock, and the question it raises}\label{sec:shock}

The OpenAI proofs of singularity formation for Euler and for Navier--Stokes have sent a shock wave
through the mathematical community, and the question being asked is plain enough. If a machine can
produce proofs of theorems better than people can, what is left for those whose main occupation is
proving theorems?

The question is legitimate and we do not propose to soften it. But two things about the present case
bear on it, and both follow from the argument above.

\emph{What was automated was the derivation, not the judgement.} The system worked on a problem
already posed, already named, already agreed to be worth settling, and already being attacked by
others along related lines~\cite{buckmaster}. It did not decide that finite-time blowup for Euler
was the question to ask. It did not decide that Euler's equations were the right model of a fluid,
which is a judgement about the world and not about symbols. And it did not connect its theorem to
anything: the step from a statement about a model to a statement about fluids is not a proof step,
and nothing in the output performs it. Sections~\ref{sec:openai} and~\ref{sec:levels} are an
instance of that gap rather than a complaint about it.

\emph{And this has happened before, to the discipline in this paper.} Machines outran hand
calculation in the 1950s, and numerical analysis did not end. It changed what it was about. The
object of study moved from finding a solution --- which the machine would do --- to bounding the
error in what the machine produced, to deciding what was worth computing, and to establishing when a
computed quantity is determined by the equations at all and when it is not. That is a harder subject
than the one it replaced, and it is the subject that produced the error control of
Section~\ref{sec:digital}.

We would therefore expect the answer for mathematics to have the same shape, without pretending to
know it. Deriving will be automated, and probably well. What will not be is choosing the question,
judging whether a model deserves the attention, supplying the instance which lets a theorem touch
something, and saying what a result means. On the evidence of this episode those are separate
activities from derivation, because here they were left undone.

\paragraph{And there is relevant symbolic mathematics, which is not theorem-proving.} If the proofs
discussed here concern the model and not fluids, that is no reflection on symbolic mathematics as
such. It is a reflection on one use of it. The most consequential symbolic work in physics has never
been the proving of theorems about given equations. It has been the \emph{formulation of the
equations}, from principles.

Equations \eqref{eq:mom}--\eqref{eq:slip} are themselves such a piece of work. Newton's second law is
a statement about particles; writing it for a continuum, deciding that the force of the surrounding
fluid is a pressure gradient, and adjoining incompressibility as a constraint with the pressure as
its multiplier, is symbolic mathematics of the first order --- and every step of it required deciding
what a fluid \emph{is}. The result is three lines carrying no parameters, from which everything in
this paper follows. Schr\"odinger's equation is the same kind of achievement for atoms, and the
reformulation of atomic physics as non-overlapping charge densities meeting at free interfaces, which
one of us has pursued under the name RealQM~\cite{realqm}, is an attempt at the same activity rather
than at theorem-proving: the work is in what the equations should be, not in what follows from them.

This is where symbolic mathematics does its heaviest lifting, where physical judgement enters
mathematics, and where, so far, no machine has done anything at all. A system which proves a theorem
about Euler's equations has been handed the hard part.

\section{Summary}\label{sec:summary}

The equations \eqref{eq:mom}--\eqref{eq:slip} are parameter free and have only geometry and forcing
as data. The mathematical questions asked of them --- existence, regularity, uniqueness --- have gone
unanswered in symbolic form because turbulent solutions have no structure for symbolic analysis to
exploit, this being what turbulence is rather than a gap in technique. In digital form, by Adaptive Euler, the same
questions are answered, in altered shape: solutions exist by exhibition, are not regular, are not
unique as fields, and are unique in the mean values anybody wants, with computed error bounds and
agreement with measurement.

A symbolic proof of finite-time blowup settles a long-open question about the model. Its solution
conserves energy exactly and so dissipates nothing, whereas a fluid dissipates and the dissipation is
the drag; and it exhibits no individual case, so nothing can be inserted into the equations,
inspected or measured. It is an exact statement about an object with no observables, set beside
approximate statements about observables which agree with experiment. Each has what the other lacks,
and only one of them is answerable to a fluid.

\begin{thebibliography}{99}

\bibitem{openaieuler} OpenAI, \emph{Finite time blowup for the Euler equation}, preprint, 2026,
57~pp.
\newblock \url{https://cdn.openai.com/pdf/315b36cd-ec98-4023-8342-93345194ece1/euler.pdf}

\bibitem{buckmaster} T.~Buckmaster and L.~Alp\"oge, \emph{Statement}, 2026.
\newblock \url{https://cims.nyu.edu/~tristanb/statement.pdf}

\bibitem{blowupHJ} J.~Hoffman and C.~Johnson, \emph{Blowup of incompressible Euler solutions},
BIT Numer.\ Math.\ \textbf{48} (2008) 285--307.

\bibitem{book} J.~Hoffman and C.~Johnson, \emph{Computational Turbulent Incompressible Flow},
Applied Mathematics: Body and Soul, vol.~4, Springer, 2007.

\bibitem{dalembert} J.~Hoffman and C.~Johnson, \emph{Resolution of d'Alembert's paradox},
J.\ Math.\ Fluid Mech.\ \textbf{12} (2010) 321--334.

\bibitem{flight} J.~Hoffman, J.~Jansson and C.~Johnson, \emph{New theory of flight},
J.\ Math.\ Fluid Mech.\ \textbf{18} (2016) 219--241.

\bibitem{dragcrisis} J.~Hoffman, \emph{Simulation of turbulent flow past bluff bodies on coarse
meshes using General Galerkin methods: drag crisis and turbulent Euler solutions}, Comput.\ Mech.\
\textbf{38} (2006) 390--402.

\bibitem{cs1} J.~Jansson, \emph{Adaptive Direct FEM simulation with Unicorn/FEniCS-HPC for CS1},
5th International Workshop on High Order CFD Methods; Zenodo, 2022.
\newblock \url{https://doi.org/10.5281/zenodo.6942599}

\bibitem{ifasd} J.~Jansson, \emph{Adaptive Euler: validation against NASA High Lift Prediction
Workshop~4, and preliminary tail-buffet and transonic-buffet simulations}, presented at IFASD~2026.

\bibitem{bkm} J.~T.~Beale, T.~Kato and A.~Majda, \emph{Remarks on the breakdown of smooth solutions
for the 3-D Euler equations}, Comm.\ Math.\ Phys.\ \textbf{94} (1984) 61--66.

\bibitem{onsager} L.~Onsager, \emph{Statistical hydrodynamics}, Nuovo Cimento \textbf{6},
Suppl.~2 (1949) 279--287.

\bibitem{realqm} C.~Johnson, \emph{A 3D multiphase continuum model for atoms and the periodic
table}, submitted, 2026.

\bibitem{mertens} A.~M.~Odlyzko and H.~J.~J.~te~Riele, \emph{Disproof of the Mertens conjecture},
J.\ reine angew.\ Math.\ \textbf{357} (1985) 138--160.

\end{thebibliography}

\end{document}
